\section{A specified integral operator order: exact indices in small rank}
\label{sec:order}

The order-index question can be answered once the lattice, lowering
operator, and ambient order are fixed.  This section supplies one
natural precise model and a complete elementary computation.  It
separates two meanings of lowering that have the same full rational
associative envelope but different integral orders.  No identification
with the integral lattice of a theta-lift chain is assumed.

\begin{theorem}[SL-09: indices in the ordinary-power string model]
\label{sl:order}
For \(n\ge1\), set
\[
\mathcal A_n=\prod_{i=0}^n
  \operatorname{Mat}_{n-i+1}(\mathbf Z).
\]
In factor \(i\), put \(d=n-i\), and choose basis
\(v_{i,0},\ldots,v_{i,d}\) in degrees
\(i,i+2,\ldots,2n-i\).  Define the integral operators
\[
L v_{i,r}=v_{i,r+1},\qquad
D v_{i,r}=r(d-r+1)v_{i,r-1},
\]
with zero outside the indicated ranges.  Let \(\pi_j\) be the
projector to degree \(j\), and let
\[
\mathcal O_n=\mathbf Z[L,D,\pi_0,\ldots,\pi_{2n}]
\subseteq\mathcal A_n
\]
be the unital associative subring.  Then
\[
[\mathcal A_1:\mathcal O_1]=1,\qquad
[\mathcal A_2:\mathcal O_2]=2^5=32,\qquad
[\mathcal A_3:\mathcal O_3]=2^{13}3^{10}=483729408.
\]
If \(D\) is replaced by the normalized unit lowering
\(\Lambda_N v_{i,r}=v_{i,r-1}\), the generated order is
\(\mathcal A_n\) for every \(n\), and its index is \(1\).
\end{theorem}

We give a direct description of the two nontrivial orders.  Matrix
indices within a string start at \(0\); \(E_{rs}\) means row \(r\),
column \(s\).  Thus \(L\) has its nonzero entries just below the
diagonal, and \(D\) has its nonzero entries just above it.

\subsection{An elementary divisibility order}

For a string of length \(d+1\), put
\[
c_r=r(d-r+1)\quad(1\le r\le d),\qquad
u_0=1,\quad u_r=\prod_{a=1}^r c_a.
\]
Define
\[
T_d=\left\{(a_{rs})\in\operatorname{Mat}_{d+1}(\mathbf Z):
       \frac{u_s}{u_r}\mid a_{rs}\text{ whenever }r<s\right\}.
\]
There is no divisibility condition on or below the diagonal.

\begin{lemma}
\label{sl:triangular-order}
The set \(T_d\) is a subring.  It is generated by \(L,D\) and all
the individual diagonal projectors of this single string.  Its
index in the full matrix ring is
\[
\prod_{0\le r<s\le d}\frac{u_s}{u_r}.
\]
\end{lemma}

\begin{proof}
For \(r<s\), consider a term \(a_{rk}b_{ks}\) in the
\((r,s)\)-entry of a product.  If \(r\le k\le s\), the two required
divisibilities supply the product
\((u_k/u_r)(u_s/u_k)=u_s/u_r\).  If \(k<r\), the second factor is
divisible by \(u_s/u_k\), itself a multiple of \(u_s/u_r\).
If \(k>s\), the first factor is divisible by \(u_k/u_r\), again
a multiple of \(u_s/u_r\).  Thus every summand meets the required
condition; \(T_d\) is a subring.

With individual diagonal projectors available, powers of \(L\)
give \(E_{rs}\) for \(r\ge s\), and powers of \(D\) give
\((u_s/u_r)E_{rs}\) for \(r<s\).  These form a
\(\mathbf Z\)-basis of \(T_d\).  The displayed product is therefore
the index of this diagonal coordinate sublattice in the full
matrix ring.
\end{proof}

The degree projectors in the theorem are not generally the
individual diagonal projectors: different strings can contribute
to the same degree.  The resulting congruence between factors
accounts for the remaining index.

\subsection{The cases n=1 and n=2}

The sums of the even and odd degree projectors are central
idempotents in the ambient ring and belong to \(\mathcal O_n\).
Consequently \(\mathcal O_n\) splits into its two parity factors.

For \(n=1\), the length-two string occupies degrees \(0,2\);
the length-one string occupies degree \(1\).  These factors are
separated by degree projectors.  The lowering coefficient on
the length-two string is \(1\).  Hence both factors are full
matrix rings, and the index is \(1\).

For \(n=2\), the odd factor is a single length-two string and
is again the full \(\operatorname{Mat}_2(\mathbf Z)\).
The even factor consists of a length-three string in degrees
\(0,2,4\) and a length-one string in degree \(2\).
For the first string,
\[
(c_1,c_2)=(2,2),\qquad (u_0,u_1,u_2)=(1,2,4).
\]
Thus \(T_2\) imposes
\[
2\mid a_{01},\qquad 2\mid a_{12},\qquad 4\mid a_{02},
\]
and has index \(2\cdot2\cdot4=2^4\) in
\(\operatorname{Mat}_3(\mathbf Z)\).

The map
\[
\rho_2:T_2\longrightarrow\mathbf F_2,\qquad
A\longmapsto a_{11}\bmod2
\]
is a surjective ring homomorphism.  Indeed in
\((AB)_{11}\), the two terms through indices \(0\) and \(2\)
are even by the displayed divisibilities.  The even order is
exactly
\[
T_2\times_{\mathbf F_2}\mathbf Z
=\{(A,z):a_{11}\equiv z\pmod2\}.
\tag{SL-09a}
\]

To prove equality, first observe that all of \(L,D,\pi_j\)
satisfy the description, so the generated ring is contained in
the right side.  Conversely, the endpoint projectors give
\((E_{00},0)\) and \((E_{22},0)\).  Multiplying \(L\) and \(D\)
by these and the middle degree projector gives
\[
(E_{10},0),\quad(E_{21},0),\quad
(2E_{01},0),\quad(2E_{12},0).
\]
Their products give \((E_{20},0)\), \((4E_{02},0)\), and
\((2E_{11},0)\).  The remaining degree projector is
\((E_{11},1)\).  These elements span every pair satisfying
\((\mathrm{SL\text{-}09a})\): all the off-diagonal coordinates
have their minimal required divisibility, and every pair of
middle diagonal entries with the same residue modulo \(2\)
is generated by \((E_{11},1)\) and \((2E_{11},0)\).
This proves equality.

Equality of the two residues is one independent index-two
condition on \(T_2\times\mathbf Z\).  Therefore the even factor
has index \(2^4\cdot2=2^5\), and so does \(\mathcal O_2\).

\subsection{The case n=3}

The odd factor is the length-three string in degrees \(1,3,5\)
together with the length-one string in degree \(3\).  It is
exactly the order just computed, with index \(2^5\).

The even factor consists of the length-four string in degrees
\(0,2,4,6\) and the length-two string in degrees \(2,4\).
Use \(E_{rs}\) for the first factor and \(e_{rs}\) for the second.
For the first factor,
\[
(c_1,c_2,c_3)=(3,4,3),\qquad
(u_0,u_1,u_2,u_3)=(1,3,12,36).
\]
The order \(T_3\) has index
\[
3\cdot12\cdot36\cdot4\cdot12\cdot3
=2^8 3^6
\tag{SL-09b}
\]
in \(\operatorname{Mat}_4(\mathbf Z)\).

Let
\[
\rho_3:T_3\longrightarrow\operatorname{Mat}_2(\mathbf F_3)
\]
be reduction modulo \(3\) of the submatrix on rows and columns
\(1,2\).  This is a ring homomorphism.  In a product of two
matrices, contributions to this submatrix through index \(0\)
contain either \(b_{01}\) or \(b_{02}\), both divisible by \(3\).
Contributions through index \(3\) contain either \(a_{13}\) or
\(a_{23}\), likewise divisible by \(3\).  The remaining terms
are precisely the product of the two middle submatrices.
The map is surjective: the middle diagonal entries and the
lower off-diagonal entry are arbitrary, and the upper
off-diagonal entry is a multiple of \(4\), which is a unit
modulo \(3\).

We claim that the even order is
\[
T_3\times_{\operatorname{Mat}_2(\mathbf F_3)}
                 \operatorname{Mat}_2(\mathbf Z).
\tag{SL-09c}
\]
Here the second map is ordinary reduction modulo \(3\).
Every generator satisfies this condition.  For the middle
link, the big-string lowering coefficient is \(4\), while
the small-string coefficient is \(1\); these coincide
modulo \(3\).  The raising coefficients are both \(1\).
The two middle degree projectors have the same matching
diagonal entries.  Endpoint terms disappear under \(\rho_3\).
This proves one containment.

For the reverse containment, endpoint projections and powers
of \(L\) supply the following isolated big-string units:
\[
(E_{00},0),\ (E_{33},0),\
(E_{10},0),\ (E_{20},0),\ (E_{30},0),\
(E_{31},0),\ (E_{32},0).
\]
For example,
\(\pi_6L^2\pi_2=(E_{31},0)\); the small string has no
degree-six component.  Powers of \(D\) similarly supply
\[
(3E_{01},0),\ (12E_{02},0),\ (36E_{03},0),\
(12E_{13},0),\ (3E_{23},0).
\]
These are all coordinates outside the middle \(2\)-by-\(2\)
block, with exactly the minimal divisibilities required
by \(T_3\).

Their products also supply the four middle kernel elements
\[
(3E_{11},0),\quad(3E_{22},0),\quad
(3E_{21},0),\quad(12E_{12},0).
\]
Explicit products, ignoring the zero second components, are
\[
E_{10}(3E_{01}),\qquad
(3E_{23})E_{32},\qquad
E_{20}(3E_{01}),\qquad
E_{10}(12E_{02}),
\]
respectively.  Finally the degree projectors and the middle
raising and lowering links provide four coupled lifts:
\[
(E_{11},e_{00}),\quad (E_{22},e_{11}),\quad
(E_{21},e_{10}),\quad (4E_{12},e_{01}).
\]
Given a pair in \((\mathrm{SL\text{-}09c})\), subtract a
linear combination of these four coupled lifts to kill its
entire small-string component.  The remaining middle
big-string entries are divisible by \(3,3,3,12\), respectively:
the upper entry is already divisible by \(4\), and its
residue modulo \(3\) now vanishes.  They are therefore a
linear combination of the four kernel elements just produced.
The isolated outer coordinates handle the rest.  This proves
\((\mathrm{SL\text{-}09c})\).

Since both residue maps are onto
\(\operatorname{Mat}_2(\mathbf F_3)\), the fibre product has
additive index
\[
\left|\operatorname{Mat}_2(\mathbf F_3)\right|=3^4
\]
in \(T_3\times\operatorname{Mat}_2(\mathbf Z)\).
Combining this with \((\mathrm{SL\text{-}09b})\) and the
odd factor yields
\[
[\mathcal A_3:\mathcal O_3]
=2^5\cdot(2^8 3^6)\cdot3^4
=2^{13}3^{10}.
\]
This completes the non-normalized part of
Theorem~\ref{sl:order}.

\subsection{Normalized lowering and the scope of H8(iv)}

For unit lowering \(\Lambda_N\), the full matrix-unit
generation was proved in SL-08, Section~\ref{sl:universal}.
For completeness, put \(B=1-L\Lambda_N\), which projects
to the bottom of each string.  Then
\[
Q_d
=B\bigl(\Lambda_N^dL^d-\Lambda_N^{d+1}L^{d+1}\bigr)B
\]
is the bottom projector of the unique string of length
\(d+1\).  Its matrix units are
\[
L^\beta Q_d\Lambda_N^\alpha,\qquad
0\le\alpha,\beta\le d.
\]
They generate the full product \(\mathcal A_n\) integrally,
even without including degree projectors as additional
generators.  This proves the final assertion of
Theorem~\ref{sl:order}.

Consequently the index of an integral operator order is
not an invariant of the full rational associative envelope
or the raising operator alone.  In the same specified
ambient order, with the same \(L\) and the same degree
projectors, the two lowering conventions give indices
\(32\) and \(1\) already for \(n=2\).

This refutes the automatic, normalization-independent
inference asserted in H8(iv), that the operator-order
index is supplied by the same binomials as the individual
exterior-algebra maps.  In particular, all divided-power
hard-Lefschetz maps on the symplectic exterior algebra for
\(n=2\) have index \(1\), whereas the explicitly specified
ordinary-power string order here has index \(32\).
These are different lattice constructions, and their
indices cannot be identified without a map respecting
the lattices and the exact operators.

The phrase \emph{with or without divided powers} also needs
a new integrality specification in this ambient order:
on a unit-step string of length at least three,
\(L^2/2!\) has a matrix entry \(1/2\), so it does not
belong to \(\mathcal A_n\).  The theorem does not assign
an index in \(\mathcal A_n\) to that noncontained ring.
It also makes no assertion that \(\mathcal O_n\) is the
image of the integral symplectic exterior lattice or the
integral lattice of an Arthur parameter.  It settles this
particular precise order model and exhibits the
normalization dependence that a further comparison must
resolve.

\subsection{Exact independent computation}

\begin{center}
\begin{tabular}{rrrr}
\toprule
\(n\)&\(\operatorname{rank}_{\mathbf Z}\mathcal A_n\)&
\([\,\mathcal A_n:\mathcal O_n\,]\)&
Index for \(\Lambda_N\)\\
\midrule
1&5&1&1\\
2&14&32&1\\
3&30&483729408&1\\
\bottomrule
\end{tabular}
\end{center}

The companion program
\texttt{normalized\_order\_indices.py} independently computes
these indices by exact column Hermite normal form.  It
vectorizes each ambient matrix unit, begins with the
\(\mathbf Z\)-span of \(1\), and repeatedly replaces a
basis matrix \(B\) by the column HNF of
\[
\bigl[B,\ R_LB,\ R_DB,\ R_{\pi_0}B,\ldots,
                         R_{\pi_{2n}}B\bigr],
\]
where \(R_g\) is right multiplication by \(g\).
At each stage this lattice is precisely the span of the
words reached by adjoining one more generator, together
with all previously reached words.  Equality with the
previous HNF therefore proves stability under every
generator, and hence equality with the generated
unital order.  The stable matrices have full rank;
their determinant absolute values are the indices.

For \(n=1,2,3\), stabilization was detected on iterations
\(2,3,4\), respectively, for both lowering conventions.
The script separately verifies, with exact rational
arithmetic, that each final lattice contains \(1\) and
all generators and that \(B^{-1}R_gB\) has integer
entries for every generator.  No floating point
determinant or probabilistic lattice test is used.
The elementary characterization and the exact HNF
computation give separate checks of the indices.

The table is supplied as
\texttt{normalized\_order\_indices.tsv}; all nonzero
coordinates of the HNF basis vectors are supplied as
\texttt{normalized\_order\_hnf\_basis.tsv}.  An additional
JSON certificate records the complete matrices and
closure histories.  These files are computation
certificates for this stated model; they are not a
formal proof-assistant certificate or a comparison
theorem for an unspecified automorphic realization.

